\abstract[Abstract]{\textbf{Background:} Fed-ISIC2019 is a widely used benchmark for federated skin
lesion classification. Its split assigns images rather than lesions to training and test sets, so images
of one lesion can appear on both sides, and the reliability of probabilities and prediction sets at
individual sites is rarely reported.
\textbf{Methods:} We audited the benchmark's six acquisition-source clients (23,247 dermoscopy images)
with image-level and lesion-grouped splits (three to five seeds per experiment), three encoder
configurations (fine-tuned and frozen EfficientNet-B0, frozen PanDerm) and up to eight federated,
personalized, centralized and local methods. We measured the apparent accuracy gain on leaked test images,
ran a randomized sibling-exposure experiment on fixed test images, and evaluated method rankings,
calibration and federated conformal coverage per client, on held-out clients and on smartphone images
(PAD-UFES-20). We also evaluated EB-ShrinkFCP, an empirical-Bayes rule for per-client conformal
thresholds. Confirmatory tests followed a version-controlled protocol fixed before the corresponding runs.
\textbf{Results:} In the official split, 0.1--93\% of each client's test images had a sibling image in
training. For FedAvg, balanced accuracy on the full image-level test set exceeded a standardized
new-lesion estimate by 8.1 points (95\% CI 4.7--11.2), and by up to 16.2 points for centralized
training. In a randomized crossover on fixed test images, having a sibling image in training raised balanced accuracy by 9.1 points (5.5--12.3) for FedAvg and by up to 28.8 points for local training. With frozen foundation-model features, method rankings were unrelated across split
types (Kendall's $\tau_b=0.05$). Calibration varied more between clients than between methods, and an apparent ECE gap between leaked and new lesions disappeared after matching sample size and composition. A global conformal threshold left some clients at 0.72--0.83
mean coverage (target 0.90). EB-ShrinkFCP kept mean client coverage near 0.90 and met its pre-specified
criterion for one of two encoders (worst-client 10th-percentile coverage +2.3 points, 0.8--3.9), so the
joint criterion was not met; a fixed shrinkage rule failed. On smartphone images, transferred coverage
fell to 0.60--0.67.
\textbf{Conclusions:} Federated dermatology benchmarks should use lesion-grouped splits and report
calibration and conformal coverage per client, with the coverage target stated. Splits, code, predictions
and the protocol history are released.}

\keywords{federated learning, data leakage, skin lesion classification, calibration, conformal prediction, benchmarking}
